\documentclass[11pt]{article}

% ============================================================
% Paper VIII: Maxwell from 4D (localized 5D Maxwell sector)
% Skeleton LaTeX (to be filled section-by-section)
% ============================================================

% -----------------------------
% Packages
% -----------------------------
\usepackage[a4paper,margin=1in]{geometry}
\usepackage[T1]{fontenc}
\usepackage{lmodern}
\usepackage{microtype}

\usepackage{amsmath,amssymb,amsthm,mathtools}
\usepackage{graphicx}
\usepackage{booktabs}
\usepackage{hyperref}
\usepackage[nameinlink,noabbrev]{cleveref}

% Optional utilities (uncomment if desired)
% \usepackage{physics}      % \dv, \pdv, \qty, etc.
% \usepackage{bm}           % bold math
% \usepackage{enumitem}     % compact lists

% -----------------------------
% Hyperref
% -----------------------------
\hypersetup{
  colorlinks=true,
  linkcolor=blue,
  citecolor=blue,
  urlcolor=blue
}

% -----------------------------
% Macros / notation
% -----------------------------
\newcommand{\dd}{\mathrm d}
\newcommand{\ii}{\mathrm i}
\newcommand{\ee}{\mathrm e}

\newcommand{\R}{\mathbb R}

% Coordinates / vectors
\newcommand{\X}{\mathbf X}   % bulk spatial (x,y,z,w)
\newcommand{\x}{\mathbf x}   % brane spatial (x,y,z)

% Fields
\newcommand{\A}{A}
\newcommand{\F}{F}
\newcommand{\J}{J}
\newcommand{\Z}{Z}

% Profile parameter (Gaussian localization width)
\newcommand{\lam}{\lambda}      % confinement width
\newcommand{\Zint}{Z_{\rm int}} % \int Z(w) dw
\newcommand{\etaQ}{\eta_Q}
\newcommand{\estar}{e_\star}
\newcommand{\qstar}{q_\star}
\newcommand{\eeff}{e_{\rm eff}}
\newcommand{\qeff}{q_{\rm eff}}
\newcommand{\Jpsi}{J_\psi}
\newcommand{\Jext}{J_{\rm ext}}
\newcommand{\Jtot}{J_{\rm tot}}

% Effective coupling
\newcommand{\muo}{\mu_0}
\newcommand{\mueff}{\mu_0^{\rm eff}}

% Gauge-fixing parameter
\newcommand{\xigf}{\xi}

% Theorem environments (optional)
\newtheorem{claim}{Claim}
\newtheorem{remark}{Remark}

% -----------------------------
% Title metadata
% -----------------------------
\title{Maxwell from the Unified 4D Toy Model:\\
Localized 5D Gauge Dynamics, Brane Reduction, and KK Corrections}
\author{Trevor Norris}
\date{\today}

\begin{document}
\maketitle

\begin{abstract}
We derive brane-effective electromagnetism from a localized \(4+1\) Maxwell sector
embedded in the unified toy model. The gauge field is a bulk potential \(A_M\)
with kinetic term \(-Z(w)F_{MN}F^{MN}/(4\mu_0)\), where the transverse profile
\(Z(w)\) localizes gauge dynamics near the brane at \(w=0\). Varying the action
yields the localized Maxwell equations
\(\partial_M(ZF^{MN})+\xi^{-1}\partial^N(\partial\!\cdot\!A)=\mu_0J^N\), the
Bianchi identities, and an exact divergence consistency identity linking the
gauge condition to current conservation. For brane-localized conserved sources,
axial gauge, and a brane-dominant zero mode, integrating over \(w\) reduces the
theory to standard \(3+1\) Maxwell with an effective coupling
\(\mu_0^{\mathrm{eff}}=\mu_0/Z_{\mathrm{int}}\), where
\(Z_{\mathrm{int}}=\int_{-\infty}^{\infty}Z(w)\,dw\). For a Gaussian profile
\(Z(w)=e^{-w^2/\lambda^2}\), \(\mu_0^{\mathrm{eff}}=\mu_0/(\lambda\sqrt{\pi})\) and
the transverse Sturm--Liouville problem yields a discrete KK tower with masses
\(m_n^2=2n/\lambda^2\) and a brane parity selection rule. The resulting static
brane potential is Coulomb plus Yukawa-suppressed corrections with a fixed
coefficient pattern; the leading deviation scales as
\(\tfrac12 e^{-2r/\lambda}\). For time-dependent sources, the brane-to-brane
response depends only on the Lorentz scalar \(k^2\) (with the retarded
prescription), ensuring brane Lorentz covariance; massive modes add causal tail
terms inside the light cone. Finally, a minimally coupled brane matter model
yields a conserved Noether current and an off-shell identity that supplies the
Maxwell source consistency condition. In the updated charge ontology, a fixed
defect branch carries \(\qstar=\etaQ\estar\) with \(\etaQ=\pm1\), while canonical
normalization of the reduced zero mode gives
\(\qeff=\qstar/\sqrt{Z_{\mathrm{int}}}\); circulation is not used to define
electric charge, and the zero-mode Maxwell limit suppresses the mixed
\((A_w,J^w,F_{\mu w})\) sector only as a controlled far-field approximation. All
principal identities and reductions are backed by a compact Wolfram Language
referee suite.
\end{abstract}

% ============================================================
\section{Introduction}
\label{sec:intro}

In the unified toy model developed across the preceding papers, the vacuum is
treated as a dynamical medium with a brane/bulk decomposition, and matter-like
excitations are associated with localized defect structures. Paper~VII introduced
a unified \(4+1\) (``5D'') formulation that cleanly separates \emph{exact} bulk
field equations from \emph{controlled} brane-effective reductions. The present
paper isolates and completes the electromagnetic sector of that framework: we
show that standard \(3+1\) Maxwell theory on the brane arises as a controlled
long-distance/low-frequency limit of a localized \(4+1\) Maxwell action, and we
quantify the leading, falsifiable deviations when transverse excitations are
relevant.

Earlier in the series, electromagnetism was developed through an ``EM-from-flow''
dictionary: with appropriate variable identifications, the homogeneous Maxwell
equations can be realized as kinematic identities and the Lorentz-force structure
can be motivated from a gauged phase-gradient (Madelung) description of matter.
While that viewpoint is useful for building intuition and for organizing
brane-observed fields, it leaves open a question that referees naturally insist
upon: \emph{where do the inhomogeneous Maxwell equations and the correct
propagator structure come from, without postulating them?} The purpose of this
paper is to supply that missing dynamical step within the unified \(4+1\) model.

In the charge ontology adopted after Paper~VII, this paper no longer uses
circulation, throat radius, or breathing variables to define electric charge.
Instead, the electric branch sign is a fixed topological puncture orientation
\(\etaQ=\pm1\), the microscopic defect-branch coupling is
\(\qstar\equiv\etaQ\estar\) with \(\estar>0\), and the observable brane strength
is controlled by the localization thickness through
\(\Zint=\int Z(w)\,\dd w\) (equivalently by canonical zero-mode normalization,
\(\qeff=\qstar/\sqrt{\Zint}\)). The familiar zero-mode Maxwell reduction is
therefore a controlled far-field brane law: it suppresses the mixed core channels
\((A_w,J^w,F_{\mu w})\) without erasing them from the microscopic ontology.

The key ingredient is a genuine \(4+1\) gauge potential \(A_M\) with a localized
Maxwell kinetic term,
\[
-\frac{Z(w)}{4\mu_0}F_{MN}F^{MN},
\]
where the nonnegative profile \(Z(w)\) confines the gauge dynamics near the brane
at \(w=0\). For concreteness (and analytic control) we adopt a Gaussian profile
\(Z(w)=\exp(-w^2/\lambda^2)\), but the logical structure applies to any profile
with finite integral \(Z_{\rm int}=\int Z(w)\,dw\). We show:

\begin{itemize}
\item \textbf{Bulk dynamics (exact):} Varying the localized action yields the
localized \(4+1\) Maxwell equations
\(\partial_M(ZF^{MN})+(1/\xi)\partial^N(\partial\!\cdot\!A)=\mu_0 J^N\), together
with exact Bianchi identities and a divergence consistency identity linking the
gauge condition to current conservation.

\item \textbf{Brane Maxwell limit (controlled):} For brane-localized sources,
axial gauge, and a brane-dominant zero mode, integrating over the transverse
direction produces standard \(3+1\) Maxwell theory with an effective coupling
\[
\mu_0^{\rm eff}=\frac{\mu_0}{Z_{\rm int}}
\quad\Rightarrow\quad
\mu_0^{\rm eff}=\frac{\mu_0}{\lambda\sqrt{\pi}}
\ \ \text{for Gaussian } Z.
\]
The same integrated equation makes explicit the correction terms that appear when
\(w\)-dependence (KK excitations) is not negligible.

\item \textbf{KK corrections (derived, quantitative):} For Gaussian localization,
the transverse Sturm--Liouville problem can be solved in closed form. The
resulting discrete KK tower yields Yukawa-suppressed corrections to Coulomb's law
with a fixed coefficient pattern; the leading correction scales as
\(\tfrac12 e^{-2r/\lambda}\) relative to the Coulomb term.

\item \textbf{Moving sources (covariant, causal):} The effective brane response is
a function of the brane Lorentz scalar \(k^2\) (with the standard retarded
prescription), implying brane Lorentz covariance for moving sources. Mode-by-mode
retarded Green functions make the causal structure explicit: massive modes add
tail-type response inside the light cone, while remaining fully causal.

\item \textbf{Source consistency (derived, not assumed):} A minimally coupled
brane matter model has a \(U(1)\) phase symmetry whose Noether current provides a
conserved source. This directly supplies the conservation condition required by
the Maxwell divergence identity, closing the logical loop between gauge symmetry,
matter dynamics, and Maxwell consistency.
\end{itemize}

\paragraph{Exact vs controlled vs regime statements.}
To keep the logical status of each step unambiguous, we adopt the following
taxonomy throughout:
(i) \emph{exact} statements are identities or Euler--Lagrange equations obtained
directly from the localized action (e.g.\ Bianchi identities, the divergence
identity);
(ii) \emph{controlled} statements require explicit reduction assumptions (e.g.\
axial gauge plus a brane-dominant zero mode); and
(iii) \emph{regime} statements interpret the controlled limit in terms of
observable scales (e.g.\ \(r\gg\lambda\), \(\omega\ll 2/\lambda\)) and summarize
the size of corrections.

\paragraph{Reproducibility and referee suite.}
All core equalities and identities in this paper are accompanied by a surgical
Wolfram Language verification suite that checks (a) the Euler--Lagrange equations
component-by-component, (b) the Bianchi identities, (c) the divergence
consistency identity, (d) the controlled brane reduction and the effective
coupling \(\mu_0^{\rm eff}\) together with the equivalent canonical relation
\(\qeff=\qstar/\sqrt{\Zint}\), (e) the Gaussian KK spectrum and brane couplings,
(f) the \(k^2\)-only dependence of the effective brane propagator for a fixed
defect branch, and (g) the matter Noether current and off-shell identity with
\(\Jpsi^\mu=\qstar j^\mu\). A concise summary of these scripts and what each
establishes is provided in \cref{app:wl_suite}.

\paragraph{Roadmap.}
\Cref{sec:conventions_action} fixes conventions and defines the localized action.
\Cref{sec:eom_identities} derives the bulk Euler--Lagrange equations and the
consistency identities. \Cref{sec:gauge_lorentz} records the gauge-variation
structure and brane Lorentz covariance. \Cref{sec:brane_reduction} presents the
controlled brane reduction and derives \(\mu_0^{\rm eff}\).
\Cref{sec:kk_coulomb} solves the Gaussian KK problem and derives the Coulomb plus
Yukawa corrections. \Cref{sec:moving_sources} extends the discussion to moving
sources and retarded structure. \Cref{sec:matter_current} derives matter current
conservation and links it to Maxwell consistency. Finally,
\Cref{sec:discussion} summarizes the validity regime, leading deviations, and
falsifiable tests.

% ============================================================
\section{Conventions and localized 5D Maxwell action}
\label{sec:conventions_action}

This paper isolates the gauge-field sector of the unified model in a standard
\(4+1\) (``5D'') spacetime notation. Concretely, we work with one time coordinate,
three brane spatial coordinates, and one transverse bulk coordinate \(w\). This is
the same localization sector used inside the full unified paper, but here we treat
it as a self-contained dynamical system whose controlled brane limit reproduces
ordinary \(3+1\) Maxwell theory (with computable corrections).

\subsection{Coordinates, indices, and metric}
\label{sec:coords_metric}

We use bulk coordinates
\begin{equation}
  x^M = (t,x,y,z,w),
  \qquad
  M,N,\ldots \in \{0,1,2,3,4\},
\end{equation}
with the brane identified as the hypersurface \(w=0\). Greek indices
\(\mu,\nu,\ldots \in \{0,1,2,3\}\) refer to brane spacetime coordinates
\(x^\mu=(t,x,y,z)\), and Latin indices \(a,b,\ldots\in\{1,2,3\}\) refer to the
brane spatial coordinates \((x,y,z)\).

We take a flat bulk metric with ``mostly plus'' signature,
\begin{equation}
  \eta_{MN}=\mathrm{diag}(-1,1,1,1,1),
  \qquad
  \eta^{MN}=(\eta_{MN})^{-1}.
\end{equation}
Indices are raised/lowered with \(\eta_{MN}\) and \(\eta^{MN}\).
Derivatives are \(\partial_M\equiv \partial/\partial x^M\), and we define the
\(4+1\) d'Alembertian
\begin{equation}
  \Box \equiv \eta^{MN}\partial_M\partial_N
  = -\partial_t^2 + \nabla^2 + \partial_w^2,
\end{equation}
where \(\nabla\) is the ordinary gradient in the brane spatial directions
\((x,y,z)\). Throughout, we work in units where the brane light-cone speed is
unity; factors of \(c\) can be restored by \(t\mapsto ct\) in the metric and
\(\partial_t\mapsto c^{-1}\partial_t\).

\subsection{Profile and normalization}
\label{sec:profile}

The hallmark of the localization mechanism is a nonnegative profile \(\Z(w)\)
multiplying the Maxwell kinetic term. In this paper we adopt a Gaussian profile,
\begin{equation}
  \Z(w) = \exp\!\left(-\frac{w^2}{\lam^2}\right),
  \label{eq:Z_gaussian}
\end{equation}
with confinement width \(\lam>0\). Its integral,
\begin{equation}
  \Zint \equiv \int_{-\infty}^{\infty}\Z(w)\,\dd w
  = \lam\sqrt{\pi},
  \label{eq:Zint}
\end{equation}
is finite, which is the key condition ensuring a normalizable zero mode and an
effective \(3+1\) gauge coupling on the brane.

All integrations by parts in \(w\) are taken with the standard assumption that the
relevant boundary terms vanish as \(|w|\to\infty\). For Gaussian \(\Z(w)\), it is
sufficient that \(\Z(w)\F^{wN}\) and \(\Z(w)\A_M\partial_w(\cdots)\) decay rapidly
enough; this holds for the localized configurations considered below and is
explicitly verified in the symbolic checks by working with smooth test fields and
then taking controlled limits.

\subsection{Fields, field strength, and gauge fixing}
\label{sec:fields}

The bulk gauge potential is a one-form (covector field) \(\A_M(x)\). Its field
strength is
\begin{equation}
  \F_{MN} \equiv \partial_M \A_N - \partial_N \A_M,
  \qquad
  \F^{MN} \equiv \eta^{MP}\eta^{NQ}\F_{PQ}.
  \label{eq:F_def}
\end{equation}
We also use the covariant divergence
\begin{equation}
  \partial\!\cdot\!\A \equiv \partial_M \A^M,
  \qquad
  \A^M \equiv \eta^{MN}\A_N.
\end{equation}

To make the kinetic operator invertible and to streamline the consistency checks,
we include a covariant gauge-fixing term with parameter \(\xigf\). This is a
standard device: \(\xigf\) labels a family of equivalent descriptions, and all
gauge-invariant observables are \(\xigf\)-independent provided the source current
is conserved. We keep \(\xigf\) arbitrary until we specialize to Lorenz gauge
\((\partial\!\cdot\!\A=0)\) when convenient.

\subsection{Action and source bookkeeping}
\label{sec:action}

The localized Maxwell sector is defined by the action
\begin{equation}
  S_{\rm EM}[\A;\Z]
  = \int \dd t\,\dd^3x\,\dd w\;
  \left[
    -\frac{\Z(w)}{4\muo}\,\F_{MN}\F^{MN}
    -\frac{1}{2\xigf\muo}\,(\partial\!\cdot\!\A)^2
    - \A_M \J^M
  \right].
  \label{eq:action_em}
\end{equation}
The source \(\J^M(x)\) is taken to be a contravariant current (a vector field),
and the interaction is fixed to the sign convention
\(\mathcal L_{\rm int}=-\A_M\J^M\). This particular bookkeeping choice is
important: it ensures that under a gauge variation \(\delta \A_M=\partial_M\chi\),
the interaction term varies as
\(\delta\mathcal L_{\rm int}=-\partial_M(\chi \J^M)+\chi\,\partial_M\J^M\),
making explicit the equivalence between gauge invariance of the source coupling
and current conservation (up to boundary terms). We emphasize this convention here
because it was historically the dominant source of component/sign mismatches in
symbolic implementations.

When embedded into the unified model, \(\J^M\) should be read as a shorthand total
source,
\begin{equation}
  \Jtot^M = \Jpsi^M + \Jext^M.
\end{equation}
In the fully coupled action, the explicit Maxwell source term is therefore
\(-\A_M\Jext^M\), while the matter current \(\Jpsi^M\) is generated separately by
minimal coupling and must not be double-counted. In the present Maxwell-only
derivation we keep the compact placeholder \(\J^M\), but in the coupled reading it
is always understood as \(\Jtot^M\). Neutralizing backgrounds or jellium
subtractions belong in \(\Jext^0\), not in a fluctuating electric-charge
dictionary.

Finally, we note a transformation-type convention that will matter later:
\(\A_\mu\) transforms as a covector under brane Lorentz transformations, while
\(\J^\mu\) transforms as a vector. With these choices, the scalar contraction
\(\A_\mu \J^\mu\) and the invariant \(\F_{\mu\nu}\F^{\mu\nu}\) are manifestly
brane-Lorentz invariant in the localized limit, as verified in the accompanying
symbolic suite.

% ============================================================
\section{Euler--Lagrange equations and consistency identities}
\label{sec:eom_identities}

This section records the exact field equations and identities implied by the
localized Maxwell action \cref{eq:action_em}. The accompanying Wolfram Language
referee suite verifies these results both in compact index form and
component-by-component for \(N\in\{t,x,y,z,w\}\), including the sign conventions
and transformation types stated in \cref{sec:conventions_action}.

\subsection{Field equations from variation}
\label{sec:eom}

Varying \(\A_N\) in \cref{eq:action_em} gives the Euler--Lagrange equations.
Using
\begin{equation}
  \delta \F_{MN}=\partial_M(\delta\A_N)-\partial_N(\delta\A_M),
\end{equation}
the variation of the kinetic term can be written as
\begin{equation}
  \delta\!\left(-\frac{\Z}{4\muo}\F_{MN}\F^{MN}\right)
  = -\frac{\Z}{2\muo}\F^{MN}\delta\F_{MN}
  = -\frac{\Z}{\muo}\F^{MN}\partial_M(\delta\A_N),
\end{equation}
where antisymmetry of \(\F^{MN}\) was used in the last step. Integrating by parts
(and dropping boundary terms, justified in \cref{sec:profile}) gives
\begin{equation}
  \delta S_{\rm kin}
  = \int \dd^5x\;\frac{1}{\muo}\,\partial_M\!\big(\Z\,\F^{MN}\big)\,\delta\A_N.
\end{equation}
Similarly, the gauge-fixing variation is
\begin{equation}
  \delta\!\left(-\frac{1}{2\xigf\muo}(\partial\!\cdot\!\A)^2\right)
  = -\frac{1}{\xigf\muo}(\partial\!\cdot\!\A)\,\partial_M(\delta\A^M)
  \;\;\Rightarrow\;\;
  \delta S_{\rm gf}
  = \int \dd^5x\;\frac{1}{\xigf\muo}\,\partial^N(\partial\!\cdot\!\A)\,\delta\A_N,
\end{equation}
where we used \(\delta\A^M=\eta^{MN}\delta\A_N\) and integrated by parts.
Finally, the source term varies as
\begin{equation}
  \delta(-\A_M\J^M) = -\J^N\,\delta\A_N.
\end{equation}
Requiring \(\delta S_{\rm EM}=0\) for arbitrary \(\delta\A_N\) yields the localized
\(4+1\) Maxwell equations,
\begin{equation}
  \partial_M\!\left(\Z(w)\F^{MN}\right)
  + \frac{1}{\xigf}\,\partial^N(\partial\!\cdot\!\A)
  = \muo\,\J^N.
  \label{eq:5d_maxwell_eom_full}
\end{equation}
We will refer to the left-hand side as the Maxwell operator,
\begin{equation}
  \mathcal{M}^N[\A] \equiv
  \partial_M\!\left(\Z\F^{MN}\right)
  + \frac{1}{\xigf}\,\partial^N(\partial\!\cdot\!\A),
  \qquad
  \mathcal{M}^N[\A]=\muo\,\J^N.
\end{equation}

\subsection{Bianchi identities (homogeneous Maxwell)}
\label{sec:bianchi}

The definition \cref{eq:F_def} implies the Bianchi identities,
\begin{equation}
  \partial_{[L}\F_{MN]} = 0,
  \label{eq:bianchi}
\end{equation}
or explicitly,
\begin{equation}
  \partial_L\F_{MN} + \partial_M\F_{NL} + \partial_N\F_{LM} = 0.
\end{equation}
These are exact, off-shell identities (they hold for any smooth \(\A_M\)),
and constitute the homogeneous Maxwell equations upon projection to brane
observables in the controlled reduction discussed later.

\subsection{Divergence identity and current conservation}
\label{sec:div_identity}

Taking \(\partial_N\) of \cref{eq:5d_maxwell_eom_full} yields a consistency
identity that ties the gauge condition to current conservation. The key point is
that the double divergence of the antisymmetric tensor \(\Z\F^{MN}\) vanishes:
\begin{equation}
  \partial_N\partial_M(\Z\F^{MN})
  = -\partial_M\partial_N(\Z\F^{MN})
  \;\;\Rightarrow\;\;
  \partial_N\partial_M(\Z\F^{MN}) = 0,
  \label{eq:double_div_zero}
\end{equation}
where we relabeled dummy indices \(M\leftrightarrow N\) and used commutativity of
partial derivatives in flat space. Therefore,
\begin{equation}
  \partial_N\big(\mathcal{M}^N[\A]-\muo\,\J^N\big)
  =
  \frac{1}{\xigf}\,\partial_N\partial^N(\partial\!\cdot\!\A)
  - \muo\,\partial_N\J^N
  =
  \frac{1}{\xigf}\,\Box(\partial\!\cdot\!\A)
  - \muo\,\partial_N\J^N.
  \label{eq:div_identity_full}
\end{equation}
In particular, in Lorenz gauge \(\partial\!\cdot\!\A=0\) the Maxwell equations
enforce \(4+1\) current conservation,
\begin{equation}
  \partial_N\J^N = 0.
  \label{eq:current_conservation}
\end{equation}
Conversely, for a conserved source \(\partial_N\J^N=0\), the gauge-fixing term
\(\propto(\partial\!\cdot\!\A)^2\) can be viewed as selecting a representative
within the gauge orbit without affecting gauge-invariant observables. This
identity is the precise statement of ``Maxwell consistency'' in the localized
\(4+1\) setting and will later be matched to the independently derived matter
Noether conservation law in \cref{sec:matter_current}.

% ============================================================
\section{Gauge structure and brane Lorentz covariance}
\label{sec:gauge_lorentz}

The localized Maxwell sector has the usual gauge structure in the bulk, modified
only by the presence of a gauge-fixing term (included here for operator
invertibility and for clean symbolic verification). In addition, because the
localization profile \(\Z(w)\) depends only on the transverse coordinate \(w\),
the theory preserves exact \(3+1\) Lorentz symmetry along the brane directions.
This section records the key structural facts that will be used later in the
brane reduction and propagator discussion.

\subsection{Gauge variation and boundary terms}
\label{sec:gauge_variation}

Consider an infinitesimal gauge transformation generated by a scalar function
\(\chi(x)\),
\begin{equation}
  \delta_\chi \A_M = \partial_M \chi.
  \label{eq:gauge_var_A}
\end{equation}
The field strength is gauge invariant,
\begin{equation}
  \delta_\chi \F_{MN}
  = \partial_M\partial_N\chi - \partial_N\partial_M\chi
  = 0,
\end{equation}
so the localized kinetic term \(-\Z\,\F^2/(4\muo)\) is unchanged. The source
interaction varies as
\begin{equation}
  \delta_\chi \mathcal{L}_{\rm int}
  = -(\delta_\chi \A_M)\J^M
  = -(\partial_M\chi)\J^M
  = -\partial_M(\chi\J^M) + \chi\,\partial_M\J^M.
  \label{eq:gauge_var_int}
\end{equation}
Thus, up to a total derivative, gauge invariance of the source coupling is
equivalent to current conservation. In particular, if \(\partial_M\J^M=0\) and
boundary terms can be dropped, then \(\delta_\chi S_{\rm int}=0\).

The gauge-fixing term explicitly breaks gauge invariance in the standard way.
Using \(\partial\!\cdot\!\A=\partial_M\A^M\), we have
\begin{equation}
  \delta_\chi(\partial\!\cdot\!\A)
  = \partial_M\delta_\chi \A^M
  = \partial_M\partial^M\chi
  = \Box \chi,
  \label{eq:divA_shift}
\end{equation}
so \((\partial\!\cdot\!\A)^2\) is not invariant. This is expected: the
gauge-fixing term is introduced only to render the kinetic operator invertible,
and physical (gauge-invariant) observables are independent of the choice of
\(\xigf\) when the source is conserved. In our referee suite we keep the
gauge-fixing term because it makes the operator comparison and divergence
identity checks unambiguous.

\subsection{Transformation types: covector vs vector}
\label{sec:covector_vector}

A recurring source of confusion in component-based checks is the difference
between how the gauge potential \(\A_\mu\) and the current \(\J^\mu\) transform
under Lorentz transformations. In the present conventions:

\begin{itemize}
\item The potential \(\A_\mu\) is a \emph{covector} (one-form) on the brane.
\item The current \(\J^\mu\) is a \emph{vector} on the brane.
\end{itemize}

Let \(\Lambda^\mu{}_\nu\) be a \(3+1\) Lorentz transformation acting on brane
coordinates \(x'^\mu=\Lambda^\mu{}_\nu x^\nu\). Then the fields transform as
\begin{equation}
  \A'_\mu(x') = (\Lambda^{-1})^\nu{}_\mu \,\A_\nu(x),
  \qquad
  \J'^\mu(x') = \Lambda^\mu{}_\nu\,\J^\nu(x).
  \label{eq:lorentz_transform_types}
\end{equation}
With these assignments, the contraction \(\A_\mu \J^\mu\) is a scalar:
\begin{equation}
  \A'_\mu(x')\J'^\mu(x')
  =
  (\Lambda^{-1})^\nu{}_\mu \A_\nu(x)\,
  \Lambda^\mu{}_\rho \J^\rho(x)
  =
  \delta^\nu{}_\rho\,\A_\nu(x)\J^\rho(x)
  =
  \A_\mu(x)\J^\mu(x).
\end{equation}
Likewise, the brane components \(\F_{\mu\nu}\) transform as a rank-2 tensor and
the invariant \(\F_{\mu\nu}\F^{\mu\nu}\) is unchanged. These transformation-type
assignments are implemented explicitly in the symbolic referee suite and are the
reason invariance checks succeed without introducing spurious sign or index
mismatches.

\subsection{\texorpdfstring{Effective brane propagator depends only on \(k^2\)}{Effective brane propagator depends only on k squared}}
\label{sec:prop_k2}

Because \(\Z(w)\) depends only on the transverse coordinate \(w\), the localized
theory is exactly invariant under the \(3+1\) Poincar\'e group acting along the
brane. In particular, after Fourier transforming in the brane coordinates,
\(\partial_\mu \to \ii k_\mu\), the brane momentum enters only through the scalar
\begin{equation}
  k^2 \equiv \eta^{\mu\nu}k_\mu k_\nu = -\omega^2 + \mathbf{k}^2,
\end{equation}
which is invariant under brane Lorentz transformations. The KK decomposition
associated with the Sturm--Liouville problem in \(w\) (worked out explicitly for
Gaussian \(\Z\) in \cref{sec:kk_coulomb}) yields an effective brane-to-brane
propagator of the schematic form
\begin{equation}
  D_{\rm eff}(k^2)
  = \sum_{n\ge 0}\frac{c_n}{k^2 + m_n^2 + \ii\epsilon},
  \label{eq:Deff_schematic}
\end{equation}
where \(m_n\) are the KK masses and \(c_n\) are the brane couplings determined by
the mode profiles at \(w=0\). Since the right-hand side depends on the brane
four-momentum only through \(k^2\), it is manifestly brane-Lorentz covariant.
This has two important consequences:

\begin{enumerate}
\item \emph{Static limit:} \(D_{\rm eff}\) reproduces a Coulomb potential from the
massless mode plus Yukawa-suppressed corrections from massive KK modes (quantified
in \cref{sec:yukawa}).

\item \emph{Moving sources:} each KK mode propagates as a covariant massless or
massive wave equation in \(3+1\), so retarded solutions for moving defect
branches retain the correct Lorentz structure mode-by-mode; summing the tower
preserves this covariance.
\end{enumerate}

In later sections we will use \cref{eq:Deff_schematic} both as a compact
organizing principle and as a bridge between the exact 5D action and the
controlled 3+1 effective description.

% ============================================================
\section{Controlled brane reduction and \texorpdfstring{$\mu_{0,\mathrm{eff}}$}{mu0,eff}}
\label{sec:brane_reduction}

The \(4+1\) equations \cref{eq:5d_maxwell_eom_full} describe a gauge field living
in the bulk with a kinetic term localized near the brane by the profile \(\Z(w)\).
In this section we explain (i) how axial gauge cleanly removes \(\A_w\) while
preserving an ordinary \(3+1\) residual gauge freedom, (ii) how brane-localized
sources are implemented, and (iii) how the integrated equations reduce to
standard \(3+1\) Maxwell with an effective coupling
\(\mueff=\muo/\Zint\). The accompanying referee suite verifies each step
symbolically, including the consistency of the \(N=w\) component under the
reduction assumptions.

\subsection{\texorpdfstring{Axial gauge reachability and residual \(3+1\) freedom}{Axial gauge reachability and residual 3+1 freedom}}
\label{sec:axial_gauge}

Start from a general configuration \(\A_M(x^\mu,w)\). Under a gauge transformation
\(\delta\A_M=\partial_M\chi\), the transverse component transforms as
\(\A_w\mapsto \A_w+\partial_w\chi\). Provided \(\A_w\) is sufficiently regular in
\(w\), we may reach axial gauge \(\A_w=0\) by choosing
\begin{equation}
  \chi_{\rm Ax}(x^\mu,w)
  \equiv -\int_{0}^{w}\A_w(x^\mu,w')\,\dd w',
  \label{eq:chi_axial}
\end{equation}
so that \(\partial_w\chi_{\rm Ax}=-\A_w\). In axial gauge the residual gauge
freedom is characterized by gauge parameters satisfying \(\partial_w\chi=0\), i.e.
\begin{equation}
  \chi(x^\mu,w)=\chi_b(x^\mu),
\end{equation}
which is precisely the usual \(3+1\) gauge freedom on the brane.

With \(\A_w=0\), the full divergence simplifies to
\begin{equation}
  \partial\!\cdot\!\A
  = \partial_M \A^M
  = \partial_\mu \A^\mu + \partial_w \A^w
  = \partial_\mu \A^\mu,
\end{equation}
so in axial gauge one has \(\partial_M \A^M=\partial_\mu \A^\mu\). In the
controlled zero-mode reduction discussed below, the clean brane Maxwell law is
obtained by imposing Lorenz gauge before integrating over \(w\), or equivalently
by fixing gauge after passing to the reduced \(3+1\) theory.

\subsection{Brane-localized sources}
\label{sec:brane_sources}

To model a conserved source confined to the brane in the controlled far-field
Maxwell reduction, we take
\begin{equation}
  \J^\mu(x^\nu,w) = \J_b^\mu(x^\nu)\,\delta(w),
  \qquad
  \J^w(x^\nu,w)=0,
  \label{eq:brane_current_delta}
\end{equation}
where \(\J_b^\mu\) is a \(3+1\) current on the brane. The condition \(\J^w=0\) is
therefore part of the reduction ansatz used to isolate the brane Maxwell sector:
it expresses the absence of net transverse current in that limit and is the
condition required by consistency of the \(N=w\) Maxwell equation under the
axial/zero-mode reduction (see below). It should not be read as a microscopic
claim that the charged core lacks mixed-sector structure. Current conservation
\(\partial_M\J^M=0\) then reduces to the usual brane condition
\(\partial_\mu \J_b^\mu=0\).

In the coupled defect reading, \(\J_b^\mu\) may represent the total brane source
\(\Jtot{}_b^\mu=\Jpsi{}_b^\mu+\Jext{}_b^\mu\). The sign of the dynamical electric
branch is then carried separately by \(\etaQ\) through
\(\Jpsi{}_b^\mu=\qstar j_b^\mu\), while any neutralizing background resides in
\(\Jext{}_b^0\).

For later convenience we also define the \(w\)-integrated current
\begin{equation}
  \J_{\rm eff}^\mu(x) \equiv \int_{-\infty}^{\infty}\J^\mu(x,w)\,\dd w,
\end{equation}
so that \cref{eq:brane_current_delta} implies \(\J_{\rm eff}^\mu=\J_b^\mu\).

\subsection{Integrated EOM and effective coupling}
\label{sec:mu_eff}

Write the \(\nu\)-components of the Maxwell equation \cref{eq:5d_maxwell_eom_full}
as
\begin{equation}
  \partial_M\!\left(\Z \F^{M\nu}\right)
  + \frac{1}{\xigf}\,\partial^\nu(\partial\!\cdot\!\A)
  = \muo\,\J^\nu,
  \qquad \nu\in\{0,1,2,3\}.
  \label{eq:eom_nu}
\end{equation}
Integrating over \(w\) gives the exact identity
\begin{equation}
  \partial_\mu \!\left(\int \Z\,\F^{\mu\nu}\,\dd w\right)
  + \underbrace{\int \partial_w\!\left(\Z\,\F^{w\nu}\right)\dd w}_{\displaystyle [\,\Z\,\F^{w\nu}\,]_{-\infty}^{+\infty}}
  + \frac{1}{\xigf}\,\partial^\nu\!\left(\int (\partial\!\cdot\!\A)\,\dd w\right)
  = \muo\,\J_{\rm eff}^\nu(x),
  \label{eq:integrated_exact}
\end{equation}
where we used \(\Z=\Z(w)\) and commuted \(\partial_\mu\) past the \(w\)-integral.
For localized field configurations the boundary term
\([\,\Z\,\F^{w\nu}\,]_{-\infty}^{+\infty}\) vanishes. It is then natural to define
a weighted brane field strength by
\begin{equation}
  \bar{\F}^{\mu\nu}(x) \equiv \frac{1}{\Zint}\int_{-\infty}^{\infty}\Z(w)\,\F^{\mu\nu}(x,w)\,\dd w,
  \label{eq:Fbar_def}
\end{equation}
so that \cref{eq:integrated_exact} becomes
\begin{equation}
  \partial_\mu \bar{\F}^{\mu\nu}
  + \frac{1}{\xigf\,\Zint}\,\partial^\nu\!\left(\int (\partial\!\cdot\!\A)\,\dd w\right)
  = \frac{\muo}{\Zint}\,\J_{\rm eff}^\nu.
  \label{eq:brane_maxwell_exact_weighted}
\end{equation}
Equation \cref{eq:brane_maxwell_exact_weighted} is an \emph{exact} integrated
brane equation, but it is not yet standard Maxwell because \(\bar{\F}^{\mu\nu}\)
is a weighted bulk average and the gauge-fixing term retains an explicit
\(w\)-integral.

To obtain the familiar \(3+1\) Maxwell form we now take a controlled zero-mode
limit: in axial gauge \(\A_w=0\), assume the brane components are \(w\)-independent
over the localization region,
\begin{equation}
  \A_\mu(x,w) \simeq a_\mu(x),
  \qquad
  \partial_w \A_\mu \simeq 0,
  \label{eq:zero_mode_ansatz}
\end{equation}
so that \(\F^{w\nu}=-\partial^w\A^\nu\simeq 0\) and
\(\F^{\mu\nu}(x,w)\simeq f^{\mu\nu}(x)\), where
\begin{equation}
  f_{\mu\nu}(x)\equiv \partial_\mu a_\nu-\partial_\nu a_\mu.
\end{equation}
Then \cref{eq:Fbar_def} gives \(\bar{\F}^{\mu\nu}\simeq f^{\mu\nu}\). However,
because the five-dimensional gauge-fixing term is not weighted by \(\Z(w)\), the
integrated gauge-fixing contribution in \cref{eq:brane_maxwell_exact_weighted} is
not naturally proportional to \(\Zint\). The clean zero-mode reduction is
therefore obtained by imposing Lorenz gauge on the five-dimensional
representative before reduction, so that \(\partial\!\cdot\!\A=0\) and the
explicit gauge-fixing term in \cref{eq:brane_maxwell_exact_weighted} vanishes
identically. Under these assumptions, \cref{eq:brane_maxwell_exact_weighted}
reduces to
\begin{equation}
  \partial_\mu f^{\mu\nu}
  = \mueff\,\J_{\rm eff}^\nu,
  \qquad
  \mueff \equiv \frac{\muo}{\Zint}.
  \label{eq:brane_maxwell_gf}
\end{equation}
For a brane-localized current \cref{eq:brane_current_delta}, \(\J_{\rm eff}^\nu=\J_b^\nu\),
and using \cref{eq:Zint} we obtain the effective coupling
\begin{equation}
  \mueff = \frac{\muo}{\lam\sqrt{\pi}}.
  \label{eq:mu0_eff}
\end{equation}
A \(3+1\) gauge-fixed representative may then be chosen after reduction, for
example \(\partial_\mu a^\mu=0\); if one wants an explicit gauge-fixing term in
the reduced theory, it should be introduced at the \(3+1\) level rather than
inferred directly from the unweighted five-dimensional term.

\paragraph{Consistency of the \(w\)-component.}
The transverse component of \cref{eq:5d_maxwell_eom_full} reads
\begin{equation}
  \partial_M\!\left(\Z \F^{Mw}\right)
  + \frac{1}{\xigf}\,\partial^w(\partial\!\cdot\!\A)
  = \muo\,\J^w.
  \label{eq:eom_w}
\end{equation}
In axial gauge \(\A_w=0\) with the zero-mode ansatz \cref{eq:zero_mode_ansatz},
we have \(\F^{\mu w}=-\partial^w\A^\mu\simeq 0\) and \(\partial^w(\partial\!\cdot\!\A)\simeq 0\),
so \cref{eq:eom_w} reduces to \(\J^w\simeq 0\). Thus the controlled brane limit is
internally consistent precisely when there is no net transverse current, matching
the source model \cref{eq:brane_current_delta}.

\paragraph{Correction terms and regime of validity.}
Equation \cref{eq:brane_maxwell_exact_weighted} makes clear what is being assumed
in the reduction. Deviations from standard \(3+1\) Maxwell arise when:
(i) \(\A_\mu\) has significant \(w\)-dependence (excited KK modes), so that
\(\bar{\F}^{\mu\nu}\neq f^{\mu\nu}\);
(ii) \(\F^{w\nu}\) does not decay rapidly enough to drop the boundary term in
\cref{eq:integrated_exact}; or (iii) \(\J^w\neq 0\).
For Gaussian localization, the KK tower and its brane couplings can be computed
explicitly, yielding Yukawa-suppressed corrections to Coulomb and to radiation
propagation; we quantify these effects in \cref{sec:kk_coulomb} and discuss their
interpretation for moving sources in \cref{sec:moving_sources}.

\paragraph{Far-field scope of the zero-mode limit.}
The assumptions \(\A_w=0\), \(\partial_w\A_\mu\simeq 0\), and \(\J^w\simeq 0\)
define a controlled \emph{far-field brane Maxwell limit}. Near a topologically
charged defect core one generically expects nontrivial mixed-sector structure in
\((A_w,J^w,F_{\mu w})\). Those channels are precisely where the puncture
orientation \(\etaQ\) and any future mixed-twist/spin microphysics would live.
Accordingly, the reduced brane Maxwell theory derived here is the effective
far-field law seen on the brane, not a statement that the microscopic charged
defect lacks odd-\(w\) or mixed core structure.

\subsection{Canonical normalization and thickness-controlled brane charge}
\label{sec:canonical_charge}

The reduction result \(\mueff=\muo/\Zint\) is the most compact way to state the
brane Maxwell limit, but the same physics can be expressed as a
thickness-controlled charge normalization. Under the zero-mode ansatz, the
effective brane action takes the schematic form
\begin{equation}
  S_{\rm EM}^{\rm eff}
  \simeq \int \dd^4x\;
  \left[
    -\frac{\Zint}{4\muo}\,f_{\mu\nu}f^{\mu\nu}
    - a_\mu \J_b^\mu
  \right].
\end{equation}
For a fixed defect branch with localized matter current
\(\Jpsi{}_b^\mu=\qstar j_b^\mu\), canonically normalizing the brane zero mode by
\begin{equation}
  a_\mu = \frac{a_\mu^{\rm can}}{\sqrt{\Zint}}
\end{equation}
gives
\begin{equation}
  S_{\rm EM}^{\rm eff}
  \simeq \int \dd^4x\;
  \left[
    -\frac{1}{4\muo}\,f_{\mu\nu}^{\rm can}f_{\rm can}^{\mu\nu}
    - a_\mu^{\rm can}\,\frac{\J_b^\mu}{\sqrt{\Zint}}
  \right].
\end{equation}
Thus the canonically normalized brane coupling of a fixed defect branch is
\begin{equation}
  \qeff=\frac{\qstar}{\sqrt{\Zint}}=\etaQ\,\eeff,
  \qquad
  \eeff\equiv \frac{\estar}{\sqrt{\Zint}}.
  \label{eq:qeff_canonical}
\end{equation}
The microscopic branch label \(\qstar\) itself remains fixed; only the reduced
brane coupling strength scales with the localization thickness. For the Gaussian
profile \(\Z(w)=\exp(-w^2/\lam^2)\), one has
\begin{equation}
  \Zint=\sqrt{\pi}\,\lam,
  \qquad
  \eeff\propto \lam^{-1/2},
\end{equation}
so thicker localization weakens the observable brane electric coupling.

% ============================================================
\section{Gaussian KK tower and corrections to Coulomb}
\label{sec:kk_coulomb}

For a finite localization integral \(\Zint\), the localized Maxwell sector admits a
normalizable zero mode that reproduces ordinary \(3+1\) Maxwell on the brane, plus
a discrete tower of massive excitations (KK modes) whose exchange produces
Yukawa-suppressed corrections. For the Gaussian profile \(\Z(w)=e^{-w^2/\lam^2}\)
these structures can be worked out in closed form. In this section we summarize:
(i) the Sturm--Liouville eigenproblem in the transverse coordinate \(w\), (ii) the
resulting spectrum and mode normalization, (iii) the brane couplings and parity
selection rule, and (iv) the resulting static potential (Coulomb plus Yukawa
corrections). These results are verified in the accompanying Wolfram Language
suite.

\subsection{Sturm--Liouville problem and spectrum}
\label{sec:kk_spectrum}

Work in axial gauge \(\A_w=0\) and, for definiteness, in Lorenz gauge for the brane
components. In the source-free region, the brane components of
\cref{eq:5d_maxwell_eom_full} separate into a brane wave operator and a transverse
operator acting on the \(w\)-profile. This motivates a mode expansion
\begin{equation}
  \A_\mu(x,w) = \sum_{n\ge 0} a^{(n)}_\mu(x)\,f_n(w),
  \label{eq:kk_expansion}
\end{equation}
where the \(f_n(w)\) are eigenfunctions of the transverse Sturm--Liouville
operator associated with the weighted Maxwell kinetic term:
\begin{equation}
  -\frac{\dd}{\dd w}\!\left(\Z(w)\,\frac{\dd f_n}{\dd w}\right)
  = m_n^2\,\Z(w)\,f_n(w).
  \label{eq:sl_problem}
\end{equation}
The weight is \(\Z(w)\), and normalizability is with respect to the inner product\\
\(\langle f,g\rangle\equiv \int_{-\infty}^{\infty}\Z(w) f(w)g(w)\,\dd w\).

For the Gaussian \(\Z(w)=\exp(-w^2/\lam^2)\), writing \(y\equiv w/\lam\) turns
\cref{eq:sl_problem} into
\begin{equation}
  \frac{\dd^2 f_n}{\dd y^2} - 2y\,\frac{\dd f_n}{\dd y} + (m_n^2\lam^2)\,f_n = 0,
\end{equation}
which is the Hermite differential equation. A complete set of polynomial
solutions is given by the physicists' Hermite polynomials,
\begin{equation}
  f_n(w) = H_n\!\left(\frac{w}{\lam}\right),
  \qquad
  m_n^2 = \frac{2n}{\lam^2},
  \qquad n=0,1,2,\ldots
  \label{eq:hermite_spectrum}
\end{equation}
with \(n=0\) the massless mode. The massive modes correspond to \(3+1\) Proca-type
propagating degrees of freedom on the brane, with masses \(m_n\).

\subsection{Mode normalization and brane couplings}
\label{sec:kk_couplings}

The natural norm induced by the action is
\begin{equation}
  \|f_n\|^2 \equiv \int_{-\infty}^{\infty}\Z(w)\,f_n(w)^2\,\dd w.
\end{equation}
For Gaussian \(\Z\) and Hermite \(f_n\), the norm is the standard Hermite orthogonality
integral,
\begin{equation}
  \|f_n\|^2
  = \int_{-\infty}^{\infty}\exp\!\left(-\frac{w^2}{\lam^2}\right)
    H_n\!\left(\frac{w}{\lam}\right)^2 \dd w
  = \lam\sqrt{\pi}\,2^n\,n!.
  \label{eq:hermite_norm}
\end{equation}

A brane-localized current \(\J^\mu(x,w)=\J_b^\mu(x)\delta(w)\) excites each mode in
proportion to its value at the brane. A convenient ``brane coupling weight'' is
\begin{equation}
  c_n \equiv \frac{f_n(0)^2}{\|f_n\|^2},
  \label{eq:cn_def}
\end{equation}
which will appear directly in the effective brane-to-brane propagator and in the
static potential.

For Hermite polynomials, \(H_n(0)=0\) for odd \(n\). Thus,
\begin{equation}
  c_{2m+1}=0,
  \qquad m=0,1,2,\ldots
  \label{eq:odd_vanish}
\end{equation}
which is a simple parity selection rule: the profile \(\Z(w)\) is even in \(w\),
the brane sits at the symmetry point \(w=0\), and odd transverse modes do not
couple to brane sources. This centered-source decoupling does not mean that
odd-\(w\) core structure is absent microscopically; it means only that such
structure can remain hidden from leading brane sources in the centered Maxwell
reduction.

For even \(n=2m\), the closed-form value \(H_{2m}(0)=(-1)^m(2m)!/m!\) gives
\begin{equation}
  c_{2m}
  = \frac{H_{2m}(0)^2}{\lam\sqrt{\pi}\,2^{2m}(2m)!}
  = \frac{1}{\lam\sqrt{\pi}}\,
    \frac{1}{4^m}\binom{2m}{m}.
  \label{eq:cn_even_closed}
\end{equation}
In particular, the zero-mode coupling is
\begin{equation}
  c_0 = \frac{1}{\lam\sqrt{\pi}} = \frac{1}{\Zint},
\end{equation}
which reproduces \(\mueff=\muo/\Zint\) from \cref{eq:mu0_eff}.

\subsection{Coulomb law and Yukawa corrections}
\label{sec:yukawa}

The KK expansion \cref{eq:kk_expansion} implies that brane momenta enter through
the Lorentz scalar \(k^2\), and the brane-to-brane propagator takes the form
\begin{equation}
  D_{\rm eff}(k^2)
  = \sum_{n\ge 0}\frac{c_n}{k^2 + m_n^2 + \ii\epsilon},
  \qquad
  m_n^2=\frac{2n}{\lam^2}.
  \label{eq:Deff_explicit}
\end{equation}
With \(\square_4=\partial_t^2-\nabla^2\) and
\(k^2=-\omega^2+\mathbf{k}^2\), this denominator is the momentum-space partner
of the mode operator \(\square_4+m_n^2\) used below. In the static limit
(\(\omega=0\)), each contributing mode produces the standard
three-dimensional Green's function for \(\nabla^2-m_n^2\), yielding a Yukawa
potential for \(m_n>0\). For a fixed defect branch on the brane, modeled by
\(\J_b^0=\qstar\,\delta^{(3)}(\mathbf x)\), the scalar potential on the brane is
\begin{equation}
  A_0(r)
  = \frac{\muo \qstar}{4\pi r}\sum_{n\ge 0} c_n\,e^{-m_n r},
  \qquad r\equiv|\mathbf x|.
  \label{eq:A0_KK_sum}
\end{equation}
Separating the zero mode and using \(\mueff=\muo c_0\), we can write
\begin{equation}
  A_0(r)
  = \frac{\mueff \qstar}{4\pi r}
    \left[
      1 + \sum_{m\ge 1}\frac{c_{2m}}{c_0}\,e^{-m_{2m} r}
    \right],
  \qquad
  m_{2m}=\frac{2\sqrt{m}}{\lam},
  \label{eq:A0_mu_eff_form}
\end{equation}
with the coefficient ratios
\begin{equation}
  \frac{c_{2m}}{c_0}
  = \frac{1}{4^m}\binom{2m}{m}
  = 1,\;\frac12,\;\frac{3}{8},\;\frac{5}{16},\;\ldots \quad (m=0,1,2,3,\ldots).
  \label{eq:ratio_cn}
\end{equation}
Therefore the leading deviation from Coulomb at large \(r\) is
\begin{equation}
  A_0(r)
  = \frac{\mueff \qstar}{4\pi r}\left[1+\frac12\,e^{-2r/\lam}
    + \mathcal{O}\!\left(e^{-2\sqrt{2}\,r/\lam}\right)\right].
  \label{eq:leading_yukawa}
\end{equation}
The corresponding radial electric field is \(E_r(r)=-\partial_r A_0(r)\), so the
zero mode gives the usual \(1/r^2\) falloff and the KK corrections are exponentially
suppressed for \(r\gg\lam\).

\paragraph{Interpretation.}
The confinement width \(\lam\) sets the crossover scale: for brane separations
\(r\gg \lam\), the zero mode dominates and the theory is experimentally
indistinguishable from Maxwell to extremely high accuracy; for \(r\sim\lam\), the
tower becomes relevant and predicts specific, computable departures from the pure
Coulomb law. In \cref{sec:moving_sources} we extend this discussion from static
sources to moving defect branches, where each KK mode propagates with its own
covariant retarded Green's function and the full response remains brane-Lorentz
covariant. Equivalently, after canonical normalization one may rewrite the same
brane field in terms of \(\qeff\), but in either normalization the Yukawa tower
modifies the field of a fixed topological charge branch; it does not imply that
electric charge varies with circulation, throat radius, or breathing.

% ============================================================
\section{Moving sources and retarded structure}
\label{sec:moving_sources}

The preceding section established that the localized Maxwell sector reduces, for a
brane-localized source, to a massless \(3+1\) mode plus a discrete tower of massive
KK excitations with computable couplings. Here we explain how this structure
extends from static sources to fully time-dependent (moving) sources, and why the
resulting brane response is causal and brane-Lorentz covariant. The key point is
that each KK mode obeys a standard covariant wave equation in \(3+1\) with a
retarded Green's function that depends only on Lorentz invariants; summing the
tower preserves that structure.

\paragraph{Mode-by-mode brane equations.}
In axial gauge with a brane-localized conserved current
\(\J^\mu(x,w)=\J_b^\mu(x)\delta(w)\), the KK decomposition
\(\A_\mu(x,w)=\sum_n a^{(n)}_\mu(x)f_n(w)\) together with the orthogonality of the
\(f_n\) under the \(\Z(w)\)-weighted inner product yields a decoupled set of
\(3+1\) equations for the mode amplitudes \(a^{(n)}_\mu(x)\). For conserved brane
sources (so that longitudinal pieces can be gauged away in the usual way), each
mode obeys a covariant Klein--Gordon-type operator with mass \(m_n\),
\begin{equation}
  \left(\square_4 + m_n^2\right)a^{(n)}_\mu(x)
  = \muo\,c_n\,\J_{b\,\mu}(x),
  \qquad
  \square_4 \equiv \partial_t^2-\nabla^2,
  \label{eq:mode_wave_eq}
\end{equation}
where the masses \(m_n^2=2n/\lam^2\) and couplings \(c_n\) were given in
\cref{eq:hermite_spectrum,eq:cn_def,eq:cn_even_closed}. The zero mode \(n=0\) is
massless and reproduces standard Maxwell propagation on the brane (with effective
coupling \(\mueff=\muo c_0\)), while the massive modes propagate as covariant
massive vector excitations.

\paragraph{Retarded Green functions.}
Let \(G^{(m)}_{\rm ret}(x-x')\) be the retarded Green function for the scalar
operator \(\square_4+m^2\),
\begin{equation}
  (\square_4+m^2)\,G^{(m)}_{\rm ret}(x-x')=\delta^{(4)}(x-x'),
  \qquad
  G^{(m)}_{\rm ret}(x-x')=0\ \ \text{for}\ t<t'.
  \label{eq:Green_def}
\end{equation}
Then the retarded solution of \cref{eq:mode_wave_eq} is
\begin{equation}
  a^{(n)}_\mu(x)
  = \muo\,c_n \int \dd^4x'\; G^{(m_n)}_{\rm ret}(x-x')\,\J_{b\,\mu}(x').
  \label{eq:mode_retarded_solution}
\end{equation}
Equivalently, in Fourier space the retarded prescription gives
\begin{equation}
  G^{(m)}_{\rm ret}(x)
  = \int\frac{\dd^4k}{(2\pi)^4}\,
    \frac{\ee^{-\ii k\cdot x}}{k^2+m^2+\ii\epsilon\,k^0},
  \qquad
  k^2\equiv \eta^{\mu\nu}k_\mu k_\nu=-\omega^2+\mathbf{k}^2,
  \label{eq:Green_retarded_fourier}
\end{equation}
where the \(\ii\epsilon k^0\) term encodes the retarded boundary condition (and is
invariant under proper orthochronous brane Lorentz transformations). The effective
brane kernel is therefore a KK-weighted sum,
\begin{equation}
  G^{\rm eff}_{\rm ret}(x)
  \equiv \sum_{n\ge 0} c_n\,G^{(m_n)}_{\rm ret}(x),
  \qquad
  a_\mu(x)=\muo\int \dd^4x'\;G^{\rm eff}_{\rm ret}(x-x')\,\J_{b\,\mu}(x'),
  \label{eq:Geff_def}
\end{equation}
and, in momentum space,
\begin{equation}
  D^{\rm eff}_{\rm ret}(k)
  = \sum_{n\ge 0}\frac{c_n}{k^2+m_n^2+\ii\epsilon\,k^0}.
  \label{eq:Deff_retarded}
\end{equation}
This makes transparent the structural statement checked in the symbolic suite:
the brane-to-brane response depends on brane momenta only through the Lorentz
scalar \(k^2\) (together with the standard retarded prescription), so there is no
preferred-frame artifact in the brane dynamics.

\paragraph{Explicit causal support and ``tails.''}
In \(3+1\), the massless retarded Green function is supported on the light cone,
\begin{equation}
  G^{(0)}_{\rm ret}(t,r)
  = \theta(t)\,\frac{\delta(t-r)}{4\pi r},
  \qquad r\equiv|\mathbf{x}|,
\end{equation}
while the massive retarded Green function has support on and \emph{inside} the
light cone (a ``tail''):
\begin{equation}
  G^{(m)}_{\rm ret}(t,r)
  = \theta(t)\left[
      \frac{\delta(t-r)}{4\pi r}
      - \theta(t-r)\,\frac{m}{4\pi}\,
        \frac{J_1\!\left(m\sqrt{t^2-r^2}\right)}{\sqrt{t^2-r^2}}
    \right],
  \label{eq:KG_retarded}
\end{equation}
where \(J_1\) is the Bessel function of the first kind. The wavefront remains
lightlike (the \(\delta(t-r)\) term), while massive modes contribute additional
causal response inside the forward light cone. Consequently, the full KK-summed
response \cref{eq:Geff_def} is manifestly causal.

\paragraph{Moving defect branches.}
For a moving defect on a fixed charge branch \(\qstar\) moving on a brane
worldline \(z^\mu(\tau)\) with four-velocity \(u^\mu=\dd z^\mu/\dd\tau\), the
conserved current can be written as
\begin{equation}
  \J_b^\mu(x) = \qstar\int \dd\tau\;u^\mu(\tau)\,\delta^{(4)}\!\big(x-z(\tau)\big).
  \label{eq:worldline_current}
\end{equation}
Substituting \cref{eq:worldline_current} into \cref{eq:mode_retarded_solution} and
summing modes gives
\begin{equation}
  a_\mu(x)
  = \muo \qstar \sum_{n\ge 0} c_n \int \dd\tau\;
    u_\mu(\tau)\,G^{(m_n)}_{\rm ret}\!\big(x-z(\tau)\big).
  \label{eq:worldline_potential}
\end{equation}
The \(n=0\) term reproduces the usual retarded Maxwell potential (the
Liénard--Wiechert structure), while the \(n\ge 1\) modes add Yukawa-suppressed and
tail-type corrections governed by the KK masses and brane couplings. In the
far-field regime \(r\gg \lam\), these corrections are exponentially small, in
accordance with \cref{eq:leading_yukawa}. The worldline current therefore
describes a moving fixed defect branch rather than an arbitrary continuously
adjustable electric charge.

\paragraph{Brane Lorentz covariance.}
Two complementary viewpoints make the covariance manifest:

\begin{enumerate}
\item \emph{Momentum-space:} the response kernel \cref{eq:Deff_retarded} depends on
brane momenta only through \(k^2\) (a Lorentz scalar), and the retarded
prescription is preserved under proper orthochronous boosts. Thus the resulting
fields transform correctly for all time-dependent sources.

\item \emph{Position-space:} the Green functions \cref{eq:KG_retarded} depend on
separations only through \(r\), \(t\), and the invariant interval
\(t^2-r^2\), together with causal step functions. Each KK mode therefore yields a
brane-covariant retarded solution, and the KK sum preserves that property.
\end{enumerate}

This completes the ``no preferred-frame'' argument for moving sources: once the
EM sector is localized by \(\Z(w)\) without breaking brane Poincar\'e symmetry, the
brane-effective dynamics are fully Lorentz covariant, with calculable deviations
from pure Maxwell behavior controlled by the KK tower.

% ============================================================
\section{Matter current conservation and Maxwell consistency}
\label{sec:matter_current}

The localized Maxwell equations \cref{eq:5d_maxwell_eom_full} require a conserved
source in order to be compatible with the gauge condition and the divergence
identity \cref{eq:div_identity_full}. In a complete theory one should not impose
source conservation by hand: it must follow from the matter dynamics and its
symmetries. In this section we show that a minimal brane matter model with
\(U(1)\) phase symmetry produces a conserved Noether current (on shell), and that
this current is precisely the object that can be used as the brane source for the
localized Maxwell sector. This closes the logical loop:
\[
\text{(matter symmetry)} \ \Rightarrow\ \partial_\mu j^\mu=0
\ \Rightarrow\ \partial_M \Jpsi^M=0
\ \Rightarrow\ \text{Maxwell consistency}.
\]

\subsection{Minimal brane matter model and Noether current}
\label{sec:noether_current}

We consider a minimal brane field \(\psi(t,\mathbf x)\) with a standard
Schr\"odinger/Gross--Pitaevskii-type Lagrangian density and minimal gauge
coupling. A representative form (sufficient for the current-conservation identity
checked in the referee suite) is
\begin{equation}
  \mathcal L_\psi
  =
  \frac{\ii\hbar}{2}\!\left(\psi^* D_t\psi - \psi (D_t\psi)^*\right)
  -\frac{\hbar^2}{2m}\,(D_i\psi)^*(D_i\psi)
  -U(|\psi|^2),
  \label{eq:L_matter}
\end{equation}
where \(U\) is an arbitrary real function of \(|\psi|^2\) and the covariant
derivatives are
\begin{equation}
  D_t \equiv \partial_t + \frac{\ii \qstar}{\hbar}A_0,
  \qquad
  D_i \equiv \partial_i - \frac{\ii \qstar}{\hbar}A_i,
  \label{eq:covariant_derivatives}
\end{equation}
with fixed defect-branch label \(\qstar=\etaQ\estar\).

The matter sector is invariant under the local \(U(1)\) transformation
\begin{equation}
  A_0 \mapsto A_0-\partial_t\chi,
  \qquad
  A_i \mapsto A_i+\partial_i\chi,
  \qquad
  \psi \mapsto \exp\!\left(\frac{\ii \qstar}{\hbar}\chi\right)\psi,
  \label{eq:matter_gauge_transform}
\end{equation}
which ensures that \(D_\mu\psi\) transforms covariantly and \(\mathcal L_\psi\)
is gauge invariant (up to total derivatives associated with the time-derivative
term). In particular, restricting \(\chi\) to a constant recovers a global phase
symmetry, and Noether's theorem yields a conserved current.

A convenient way to extract the current (and the approach used in the symbolic
suite) is to probe a \emph{local} phase variation
\(\psi\to e^{\ii \epsilon(x)}\psi\), \(\psi^*\to e^{-\ii \epsilon(x)}\psi^*\)
and expand to first order in \(\epsilon\). One finds a current \(j^\mu\) on the
brane with time component
\begin{equation}
  j^0 = |\psi|^2,
  \label{eq:j0}
\end{equation}
and spatial components
\begin{equation}
  j^i
  = \frac{\hbar}{2mi}\Big(\psi^* D^i\psi - (D^i\psi)^*\psi\Big)
  = \frac{\hbar}{m}\,\mathrm{Im}\!\left(\psi^* D^i\psi\right),
  \label{eq:ji_covariant}
\end{equation}
which is manifestly gauge covariant. The corresponding \emph{charge current} is
\begin{equation}
  \Jpsi^\mu \equiv \qstar\,j^\mu.
  \label{eq:charge_current_def}
\end{equation}
This minimal brane matter model is used here as a representative
source-consistency closure, not as a claim that the entire ambient vacuum
condensate is uniformly charged. In the defect reading, \(\psi\) may be treated as
a localized, unit-normalized defect profile carrying fixed total charge
\(\qstar\), while any neutralizing background is encoded separately in
\(\Jext^0\). In the full coupled theory, \(\Jpsi^\mu\) is the dynamical source
generated by minimal coupling, while the explicit Maxwell action term carries only
\(\Jext^\mu\).

\subsection{Off-shell identity and linkage to Maxwell divergence}
\label{sec:off_shell_identity}

The local phase-probe method yields not only the explicit form of the Noether
current but also an \emph{off-shell} identity (Noether identity) relating its
divergence to the matter Euler--Lagrange expressions. Denote the brane
Euler--Lagrange operators by
\begin{equation}
  \mathrm{EL}_{\psi} \equiv \frac{\partial\mathcal L_\psi}{\partial \psi}
  -\partial_\mu\!\left(\frac{\partial\mathcal L_\psi}{\partial(\partial_\mu\psi)}\right),
  \qquad
  \mathrm{EL}_{\psi^*} \equiv \frac{\partial\mathcal L_\psi}{\partial \psi^*}
  -\partial_\mu\!\left(\frac{\partial\mathcal L_\psi}{\partial(\partial_\mu\psi^*)}\right).
\end{equation}
Then the phase symmetry implies the identity
\begin{equation}
  \partial_\mu j^\mu
  + \frac{\ii}{\hbar}\Big(\psi^*\mathrm{EL}_{\psi^*}-\psi\,\mathrm{EL}_{\psi}\Big)
  = 0.
  \label{eq:noether_identity}
\end{equation}
On shell, \(\mathrm{EL}_\psi=\mathrm{EL}_{\psi^*}=0\), this reduces to the usual
continuity equation
\begin{equation}
  \partial_\mu j^\mu = 0,
  \qquad\text{equivalently}\qquad
  \partial_\mu \Jpsi^\mu = 0.
  \label{eq:brane_current_conservation}
\end{equation}
The referee suite verifies \cref{eq:noether_identity} explicitly for the minimal
matter model used there (including the corrected sign convention in the identity
residual).

To embed the matter source into the localized \(4+1\) Maxwell sector, we use the
brane-localized current model
\begin{equation}
  \J^\mu(x,w) = \Jpsi^\mu(x)\,\delta(w),
  \qquad
  \J^w(x,w)=0,
  \label{eq:J_bulk_from_brane}
\end{equation}
as in \cref{sec:brane_sources}. Since \(\delta(w)\) depends only on \(w\), the
bulk divergence reduces distributionally to
\begin{equation}
  \partial_M \J^M
  = \partial_\mu\!\big(\Jpsi^\mu\delta(w)\big)
    + \partial_w(0)
  = \delta(w)\,\partial_\mu \Jpsi^\mu.
\end{equation}
Therefore, the on-shell brane conservation law \cref{eq:brane_current_conservation}
implies bulk conservation \(\partial_M\J^M=0\). This is exactly the condition
required by Maxwell consistency: inserting \(\partial_M\J^M=0\) into the
divergence identity \cref{eq:div_identity_full} shows that in Lorenz gauge
\((\partial\!\cdot\!\A=0)\) the bulk equations are mutually compatible and the
gauge-fixed operator closes without contradiction.

In summary, the same symmetry (phase invariance) that makes the matter sector
well-posed also supplies the conserved source required by the localized Maxwell
sector. In the controlled brane limit, this guarantees that the effective \(3+1\)
Maxwell equations are sourced by a current that is conserved by construction,
rather than by assumption.

% ============================================================
\section{Discussion: validity regime, deviations, and falsifiable tests}
\label{sec:discussion}

The localized \(4+1\) Maxwell sector yields ordinary \(3+1\) electromagnetism on
the brane only in a controlled regime. Outside that regime, the theory makes
specific, quantitative predictions: Yukawa-suppressed corrections to Coulomb, a
discrete set of massive response channels with threshold behavior in frequency,
and (for time-dependent sources) causal ``tail'' terms inside the light cone.
Because the correction structure is fixed by the localization profile \(\Z(w)\),
these deviations are falsifiable rather than tunable.

\subsection{\texorpdfstring{When the \(3+1\) Maxwell limit holds}{When the 3+1 Maxwell limit holds}}
The exact integrated brane equation \cref{eq:brane_maxwell_exact_weighted} shows
that the reduction to standard Maxwell is controlled by (i) the dominance of a
\(w\)-independent zero mode in the field configuration and (ii) the absence of
transverse current. A sufficient set of conditions for the Maxwell regime is:

\begin{enumerate}
\item \textbf{Zero-mode dominance:} the brane components are approximately
\(w\)-independent within the support of \(\Z(w)\),
\(\A_\mu(x,w)\simeq a_\mu(x)\) and \(\partial_w \A_\mu\simeq 0\), as in
\cref{eq:zero_mode_ansatz}. Equivalently, the typical brane length/time scales
are long compared to the confinement width:
\begin{equation}
  r \gg \lam,
  \qquad
  \omega \ll \frac{2}{\lam},
\end{equation}
so that KK excitations are either not sourced efficiently or remain evanescent.

\item \textbf{No transverse source flux:} \(\J^w=0\), so the \(N=w\) Maxwell
equation \cref{eq:eom_w} is consistent with axial gauge and the zero-mode limit.

\item \textbf{Localized boundary behavior:} \(\Z(w)\F^{w\nu}\) decays sufficiently
fast as \(|w|\to\infty\) so that the boundary term in \cref{eq:integrated_exact}
vanishes (automatic for the localized configurations considered here).

\item \textbf{Conserved brane current:} \(\partial_\mu \J_b^\mu=0\), which follows
from the matter dynamics via the Noether identity \cref{eq:noether_identity}.
\end{enumerate}

Under these conditions the theory reduces to the familiar \(3+1\) Maxwell
equation \cref{eq:brane_maxwell_gf} with effective coupling
\(\mueff=\muo/\Zint\); a convenient gauge such as \(\partial_\mu a^\mu=0\) may
then be chosen within the reduced theory, and all gauge-invariant brane
observables coincide with Maxwell to high accuracy.

In the updated charge ontology, this far-field Maxwell limit describes the brane
field of a fixed defect branch with sign \(\etaQ\) and effective brane coupling
\(\qeff=\qstar/\sqrt{\Zint}\). What is being suppressed here is the mixed core
sector \((A_w,J^w,F_{\mu w})\), not the underlying puncture orientation itself.

\subsection{Leading deviation structure}
The Gaussian profile fixes the KK spectrum and brane couplings, yielding a rigid
deviation pattern rather than an arbitrary modification. For static sources the
brane potential takes the KK-summed form \cref{eq:A0_KK_sum}; separating the zero
mode gives the leading correction \cref{eq:leading_yukawa},
\begin{equation}
  A_0(r)=\frac{\mueff \qstar}{4\pi r}
  \left[1+\frac12\,\ee^{-2r/\lam}
  +\mathcal O\!\left(\ee^{-2\sqrt2\,r/\lam}\right)\right].
\end{equation}
Two points are worth emphasizing:

\begin{itemize}
\item The \emph{range} of the leading deviation is fixed by the first brane-coupled
mass \(m_2=2/\lam\) (odd modes do not couple; \cref{eq:odd_vanish}).

\item The \emph{amplitude} of the leading Yukawa term is also fixed (here \(1/2\))
by the Gaussian mode structure (\cref{eq:ratio_cn}). In this sense the Gaussian
localization makes a parameter-light, falsifiable prediction: once \(\lam\) is
chosen, the leading correction is determined.
\end{itemize}

These Yukawa corrections modify the field profile of a fixed topological charge
branch. They do not imply any dependence of electric charge on circulation,
throat radius, or geometry breathing.

For time-dependent sources, each KK mode obeys a covariant massive wave equation
\cref{eq:mode_wave_eq} with retarded solution \cref{eq:mode_retarded_solution}.
Below the first brane-coupled threshold \(\omega<m_2\), massive modes contribute
only near-field/evanescent corrections and the response is extremely close to
Maxwell at distances \(r\gg\lam\). Above threshold, \(\omega>m_2\), additional
propagating channels open with dispersion
\begin{equation}
  \omega^2 = k^2 + m_n^2,
\end{equation}
and group velocity \(v_g = k/\omega < 1\). This predicts frequency-dependent
response (and, in principle, additional polarization content) controlled by the
discrete tower \(m_n^2=2n/\lam^2\) and couplings \(c_n\).

\subsection{Falsifiable tests and observational handles}
The localization sector can be tested in several complementary ways. The
discussion below is organized by which assumption in the controlled reduction is
being probed.

\paragraph{(i) Precision tests of Coulomb's law (static).}
The cleanest signature is a Yukawa departure with \emph{fixed coefficient pattern}
set by \(\Z(w)\). For Gaussian localization the leading correction has coefficient
\(1/2\) and range \(\lam/2\) in the exponent (\cref{eq:leading_yukawa}). Any
experimental bound on Yukawa deviations from \(1/r\) behavior can be translated
directly into a bound on \(\lam\). Conversely, a detected deviation would overconstrain
the model because higher-mode coefficients are not free: the subleading terms are
predicted to follow \cref{eq:ratio_cn} with masses \(\{2\sqrt{m}/\lam\}_{m\ge1}\).

\paragraph{(ii) Threshold behavior and dispersion (time-dependent).}
For driven sources at frequency \(\omega\), the theory predicts a qualitative
change in response when \(\omega\) crosses the first coupled KK mass \(m_2=2/\lam\).
Below threshold, KK contributions are evanescent; above threshold they propagate
with subluminal group velocity. If \(\lam\) were large enough to place \(m_2\) in
an accessible band, this would appear as a sharp onset of dispersive components or
additional attenuation channels. If no such effect is seen up to some frequency,
one obtains a lower bound on \(m_2\), hence an upper bound on \(\lam\).

\paragraph{(iii) Retarded ``tail'' structure (time-domain).}
Massive modes produce causal tails inside the forward light cone
(\cref{eq:KG_retarded}), implying that a sharp electromagnetic pulse sourced on
the brane can generate a small, model-predicted after-response beyond the
light-cone wavefront. Detecting (or constraining) such tails provides a direct
probe of massive KK content distinct from static-force measurements.

\paragraph{(iv) Charge conservation and possible leakage signatures.}
The controlled reduction requires \(\J^w=0\) and yields bulk conservation
\(\partial_M\J^M=0\) when the brane matter equations enforce
\(\partial_\mu \J_b^\mu=0\) (\cref{eq:brane_current_conservation}). Any empirical
signature of apparent charge nonconservation on the brane (or of systematic
violations of the Maxwell divergence identity) would falsify the source model
\cref{eq:brane_current_delta} and force the inclusion of transverse current or
additional degrees of freedom. Conversely, the framework offers a precise place
to parametrize such effects: \(\J^w\neq 0\) is not ``forbidden,'' but it moves one
outside the Maxwell regime and predicts accompanying modifications in the \(N=w\)
equation \cref{eq:eom_w}.

\paragraph{(v) Profile-shape tests beyond the Gaussian.}
While this paper adopts a Gaussian \(\Z(w)\) for analytic control, the general
framework depends on \(\Z(w)\) only through the Sturm--Liouville problem
\cref{eq:sl_problem} and the resulting tower \(\{m_n,c_n\}\). Different profiles
produce different discrete spectra and coupling patterns. Thus, experimental
constraints can be interpreted either as bounds on \(\lam\) \emph{within} the
Gaussian class, or more broadly as constraints on the admissible profile shapes
that yield the observed near-perfect Maxwell behavior on the brane.

\subsection{Relation to the unified model and next derivations}
Within the broader unified program, the localized Maxwell sector plays two roles:
(i) it supplies a genuine gauge dynamics on the brane (including radiation and
Lorentz-force coupling), and (ii) it provides a controlled place to quantify
deviations (KK tails, Yukawa corrections) that can be linked to the transverse
structure of the unified geometry. A natural next step is to derive or constrain
\(\Z(w)\) and \(\lam\) from the throat/defect microphysics rather than treating
them as phenomenological response data. The referee suite provided with this
paper is designed to make that future step modular: any proposed \(\Z(w)\) can be
dropped into the same symbolic pipeline to produce the corresponding
\(\mueff\), \(\qeff\), KK spectrum, and brane deviation pattern. The mixed
\((A_w,J^w,F_{\mu w})\) sector remains the natural place for future defect-core
twist/spin dynamics, but that microphysics is outside the scope of the present
Maxwell-localization paper.

% ============================================================
\section{Conclusion}
\label{sec:conclusion}

We isolated the gauge-field sector of the unified \(4+1\) framework and showed
that ordinary \(3+1\) Maxwell theory on the brane emerges as a controlled,
script-verified limit of a localized \(5\)D Maxwell action. The core ingredients
are: (i) a transverse localization profile \(\Z(w)\) multiplying the Maxwell
kinetic term, (ii) standard gauge structure (with an explicit gauge-fixing term
used here to make operator inversion and consistency checks transparent), and
(iii) a brane-localized conserved source.

Starting from the localized action \cref{eq:action_em}, variation with respect to
\(\A_N\) yields the bulk Euler--Lagrange equations
\cref{eq:5d_maxwell_eom_full}. The same framework implies the homogeneous Maxwell
equations via the Bianchi identities \cref{eq:bianchi}, and it enforces a precise
divergence consistency identity \cref{eq:div_identity_full} relating the gauge
condition to current conservation. We also made explicit the gauge-variation
structure of the source coupling \cref{eq:gauge_var_int} and the correct
transformation-type assignments (covector \(\A_\mu\), vector \(\J^\mu\)) required
for clean brane Lorentz-covariance checks.

Under axial gauge and a brane-dominant zero-mode ansatz, integrating the bulk
equations over \(w\) reduces the dynamics to standard \(3+1\) Maxwell theory with
a renormalized coupling
\begin{equation}
  \mueff = \frac{\muo}{\Zint},
  \qquad
  \Zint=\int_{-\infty}^{\infty}\Z(w)\,\dd w.
\end{equation}
For the Gaussian profile \(\Z(w)=\exp(-w^2/\lam^2)\), \(\Zint=\lam\sqrt{\pi}\) and
\(\mueff=\muo/(\lam\sqrt{\pi})\). Beyond this Maxwell regime, the same
localization mechanism predicts a discrete KK tower governed by the
Sturm--Liouville problem \cref{eq:sl_problem}. For Gaussian \(\Z\), the spectrum
is \(m_n^2=2n/\lam^2\) with an even-mode coupling selection rule
(\(c_{2m+1}=0\)), yielding a Coulomb potential plus Yukawa-suppressed corrections
whose leading term has fixed coefficient \(1/2\) and range set by \(\lam\)
(\cref{eq:leading_yukawa}). For time-dependent sources, each KK mode propagates
with a covariant retarded Green function, so the full brane response remains
causal and brane-Lorentz covariant, with calculable dispersive/tail corrections
(\cref{sec:moving_sources}).

Finally, we closed the consistency loop on the source side. A minimal brane
matter model with \(U(1)\) phase symmetry yields a Noether current \(j^\mu\) with
\(j^0=|\psi|^2\) and gauge-covariant spatial components \cref{eq:ji_covariant},
together with the off-shell identity \cref{eq:noether_identity} implying
\(\partial_\mu j^\mu=0\) on shell. Embedding the associated charge current as a
brane-localized bulk source \(\J^\mu=\Jpsi^\mu\delta(w)\), \(\J^w=0\) ensures
\(\partial_M\J^M=0\), precisely the condition required by the Maxwell divergence
identity and by gauge consistency.

In summary, the localized Maxwell sector provides a complete, falsifiable
mechanism for obtaining brane electromagnetism from the unified \(4+1\) model:
standard Maxwell dynamics emerge at long distances and low frequencies, while
well-structured deviations (Yukawa corrections, discrete thresholds, and causal
tails) appear when transverse excitations are relevant. The accompanying symbolic
referee suite was designed to make these claims reproducible and to enable direct
future substitutions of alternative localization profiles \(\Z(w)\) derived from
the underlying defect/throat microphysics. In this updated ontology, the electric
branch sign is fixed by \(\etaQ\), the observable brane charge strength is
thickness-controlled through \(\Zint\) (equivalently \(\qeff\)), and the zero-mode
reduction suppresses rather than eliminates the mixed core sector. This paper
derives the localized Maxwell dynamics and the thickness-controlled brane charge
strength; it does not derive spin or the microscopic mixed-twist defect-core
dynamics.

\section*{Acknowledgments}
This work was carried out by an independent researcher. Generative-AI tools were
used during the development of the manuscript to help translate conceptual ideas
into mathematical form, assist with symbolic manipulations, prototype and revise
companion Mathematica and SymPy verification scripts, and improve the
exposition. The conceptual direction, hypotheses, and overall structure of the
work were directed by the author. The resulting arguments and calculations were
checked for internal consistency using the accompanying symbolic scripts, but
the work should be regarded as exploratory and would benefit from review by
subject-matter experts. Any errors, omissions, or misunderstandings remain the
author's.

% ============================================================
\appendix

% ============================================================
\section{Referee verification suite summary (Wolfram Language)}
\label{app:wl_suite}

This paper is accompanied by a compact Wolfram Language (WL) verification suite
that checks the principal identities and reductions used in the main text. The
suite is designed to be referee-friendly: each script prints the exact symbolic
Lagrangian it varies (when applicable), converts \texttt{EulerEquations} output to
expressions before comparison (to avoid \texttt{Equal[\dots]} mismatches), and
tests all nontrivial identities by simplifying the \emph{residual} to zero under
the declared conventions. In particular, the suite locks in the sign conventions
and transformation types used throughout the paper: metric
\(\eta_{MN}=\mathrm{diag}(-1,1,1,1,1)\), field strength
\(\F_{MN}=\partial_M\A_N-\partial_N\A_M\), source coupling written in compact
Maxwell-only form as \(\mathcal L_{\rm int}=-\A_M\J^M\) with \(\J^M\)
contravariant and \(\A_M\) covariant, and Gaussian localization
\(\Z(w)=\exp(-w^2/\lam^2)\). In the fully coupled embedding, the explicit
Maxwell source term is instead \(-\A_M\Jext^M\), while the dynamical matter
current \(\Jpsi^M\) is generated separately by minimal coupling so that
\(\Jtot^M=\Jpsi^M+\Jext^M\).

\paragraph{(A) \texttt{maxwell\_from\_4d.wl}: action variation, EOM, Bianchi, divergence identity, brane reduction.}
This is the core referee harness. It defines the localized Maxwell action with
Gaussian \(\Z(w)\), includes the Lorenz-type gauge-fixing term, and varies with
respect to \(\A_N\). It verifies:
(i) the full \(4+1\) Euler--Lagrange equations match
\(\partial_M(\Z\F^{MN})+(1/\xigf)\partial^N(\partial\!\cdot\!\A)=\muo\J^N\) for all
components \(N\in\{t,x,y,z,w\}\) (residuals simplify to \(0\));
(ii) the Bianchi identities \(\partial_{[L}\F_{MN]}=0\) hold identically for all
index triples (residuals simplify to \(0\));
(iii) the divergence consistency identity
\(\partial_N(\mathrm{EOM}^N)=\xigf^{-1}\Box(\partial\!\cdot\!\A)-\muo\,\partial_N\J^N\)
holds exactly (residual \(0\));
and (iv) under a controlled brane/zero-mode reduction (\(\A_w=0\),
\(\partial_w\A_\mu=0\)), the integrated equations reproduce standard \(3+1\)
Maxwell with \(\mueff=\muo/\Zint=\muo/(\lam\sqrt{\pi})\). The same harness also
prints the updated charge bookkeeping
\(\qstar=\etaQ\estar\), \(\Jtot^M=\Jpsi^M+\Jext^M\), and the equivalent
canonical brane-coupling relation \(\qeff=\qstar/\sqrt{\Zint}\).

\paragraph{(B) \texttt{maxwell\_symmetry\_checks.wl}: gauge variation and brane Lorentz invariance.}
This script checks the symmetry structure emphasized in
\cref{sec:gauge_lorentz}. It verifies the gauge variation of the interaction term,
\(\delta\mathcal L_{\rm int}=-\partial_M(\chi\J^M)+\chi\,\partial_M\J^M\), and the
shift \(\delta(\partial\!\cdot\!\A)=\Box\chi\) implied by \(\delta \A_M=\partial_M\chi\).
It also verifies brane Lorentz invariance by implementing the correct
transformation types: \(\A_\mu\) as a covector and \(\J^\mu\) as a vector. In
addition, it checks the Paper~VII matter-gauge convention
\(A_0\mapsto A_0-\partial_t\chi\), \(A_i\mapsto A_i+\partial_i\chi\),
\(\psi\mapsto e^{\ii\qstar\chi/\hbar}\psi\), confirming that the corresponding
covariant derivatives transform homogeneously. With this bookkeeping, it
confirms invariance of \(\F_{\mu\nu}\F^{\mu\nu}\) and \(\A_\mu\J^\mu\) under
representative boosts (residuals simplify to \(0\)).

\paragraph{(C) \texttt{maxwell\_kk\_coulomb\_propagator.wl}: Gaussian KK spectrum, couplings,\\Coulomb+Yukawa, and \(D_{\rm eff}(k^2)\).}
This script establishes the transverse Sturm--Liouville eigenproblem
\(-\frac{\dd}{\dd w}(\Z f')=m^2\Z f\) for Gaussian \(\Z\) and verifies the Hermite
solutions \(f_n(w)=H_n(w/\lam)\) with spectrum \(m_n^2=2n/\lam^2\). It validates the
mode norms \(\|f_n\|^2=\lam\sqrt{\pi}\,2^n n!\) by integer spot-checks and then
uses the closed form to compute brane couplings
\(c_n=f_n(0)^2/\|f_n\|^2\), confirming the parity selection rule \(c_{2m+1}=0\) and
the even-mode formula \(\frac{c_{2m}}{c_0}=\frac{1}{4^m}\binom{2m}{m}\). It then
derives the Coulomb zero-mode potential and the KK Yukawa corrections for a fixed
defect branch labeled by \(\qstar\), records the equivalent canonical coupling
\(\qeff=\qstar/\sqrt{\Zint}\), and notes explicitly that odd transverse structure
can decouple from centered brane sources without being absent microscopically.
Finally, it checks that truncated Euclidean (and Minkowski) effective propagators
depend only on the brane Lorentz scalar \(k^2\).

\paragraph{(D) \texttt{maxwell\_axial\_gauge\_brane\_sources.wl}: axial gauge reachability, residual \(3+1\) gauge, \(\delta(w)\) sources, and integrated EOM.}
This script demonstrates explicit axial gauge reachability via the parameter
\(\chi_{\rm Ax}(w)=-\int_0^w \A_w\,\dd w'\), and shows that the residual gauge
freedom is \(w\)-independent (standard \(3+1\) gauge parameter). It implements a
brane-localized current \(\J^\mu=\J_b^\mu(x)\delta(w)\), \(\J^w=0\), and verifies:
(i) the \(N=w\) Maxwell equation is consistent with axial gauge and the zero-mode
ansatz; and (ii) integrating the \(N=\nu\) equations over \(w\) yields precisely
the \(3+1\) Maxwell form with \(\mueff=\muo/\Zint\). The harness also records the
interpretive caveat used in the paper: \(\J^w=0\) is a controlled far-field
Maxwell-limit condition, not a claim that the microscopic defect core lacks mixed
\((A_w,J^w,F_{\mu w})\) structure.

\paragraph{(E) \texttt{maxwell\_moving\_charges\_propagator.wl}: covariant response for time-dependent sources.}
This script provides a complementary momentum-space verification of the statements
in \cref{sec:moving_sources}. It demonstrates explicitly that \(k^2\) is brane
boost invariant and that the KK-truncated effective propagator depends only on
\(k^2\), and it supplies the corresponding Minkowski \(\ii\epsilon\) prescriptions
for retarded structure mode-by-mode. In the static limit, it reproduces the
Coulomb+Yukawa structure for a fixed defect branch labeled by \(\qstar\); more
generally it supports the statement that moving defect branches do not introduce
preferred-frame artifacts in the brane response.

\paragraph{(F) \texttt{matter\_current\_conservation.wl}: Noether current and off-shell identity.}
This script derives the brane matter Noether current by probing a local phase
variation \(\epsilon(t,x,y,z)\) in a minimally coupled matter Lagrangian of the
form summarized in \cref{eq:L_matter}. It verifies:
(i) constant phase transformations leave \(\mathcal L_\psi\) invariant;
(ii) the extracted number current satisfies \(j^0=|\psi|^2\) and the expected
Paper~VII gauge-covariant spatial form; and (iii) the off-shell Noether identity
\(\partial_\mu j^\mu + \frac{\ii}{\hbar}(\psi^*\mathrm{EL}_{\psi^*}-\psi\,\mathrm{EL}_\psi)=0\)
holds exactly (residual simplifies to \(0\)) with the fixed sign convention used
in this paper. It then records the Maxwell source as \(\Jpsi^\mu=\qstar j^\mu\),
with \(\qstar=\etaQ\estar\), matching the localized-defect / background-split
bookkeeping used in the main text.

\paragraph{Reproducibility note.}
All scripts are intended to be run in a standard WL kernel (e.g.\ via
\texttt{wolframscript -file <file.wl>}). Each prints pass/fail messages and
(where appropriate) the residual expressions before simplification. The suite is
structured so that replacing the profile \(\Z(w)\) by an alternative localization
ansatz requires changes in only one place, after which the same EOM, identity,
KK-spectral, reduction, and charge-normalization checks can be re-run verbatim.
% ============================================================

% ============================================================
\section{Sign conventions and common pitfalls}
\label{app:conventions_pitfalls}

This appendix records the sign, index, and transformation conventions used in the
paper and highlights the few implementation pitfalls that historically cause
spurious discrepancies (especially in component-based symbolic checks). The goal
is to make it easy for readers and referees to reproduce results without
re-deriving the bookkeeping.

\subsection{Frozen conventions used throughout}

\paragraph{Coordinates and metric.}
We use bulk coordinates \(x^M=(t,x,y,z,w)\) with\\
\(\eta_{MN}=\mathrm{diag}(-1,1,1,1,1)\). The bulk d'Alembertian is
\(\Box=\eta^{MN}\partial_M\partial_N=-\partial_t^2+\nabla^2+\partial_w^2\).
Brane indices are \(\mu,\nu\in\{0,1,2,3\}\) and spatial brane indices
\(i,j\in\{1,2,3\}\).

\paragraph{Field strength.}
The field strength is defined by
\begin{equation}
  \F_{MN}=\partial_M\A_N-\partial_N\A_M,
  \qquad
  \F^{MN}=\eta^{MP}\eta^{NQ}\F_{PQ}.
\end{equation}
With these definitions, the Bianchi identities
\(\partial_{[L}\F_{MN]}=0\) hold identically.

\paragraph{Localization profile.}
The localization function \(\Z(w)\) multiplies the Maxwell kinetic term. In the
main text we take the Gaussian
\(\Z(w)=\exp(-w^2/\lam^2)\), with finite integral
\(\Zint=\int\Z\,\dd w=\lam\sqrt{\pi}\).

\paragraph{Source coupling sign and index placement.}
The compact Maxwell-only source term is written as
\begin{equation}
  \mathcal L_{\rm int} = -\A_M \J^M,
\end{equation}
where \(\J^M\) is treated as \emph{contravariant} (a vector field), and \(\A_M\)
is \emph{covariant} (a one-form). In the fully coupled embedding, however, the
explicit Maxwell action carries only the external/background source,
\(-\A_M\Jext^M\), while the dynamical matter current \(\Jpsi^M\) is generated by
minimal coupling, so that \(\Jtot^M=\Jpsi^M+\Jext^M\). With the compact
convention, a gauge transformation \(\delta \A_M=\partial_M\chi\) yields
\begin{equation}
  \delta \mathcal L_{\rm int}
  = -(\partial_M\chi)\J^M
  = -\partial_M(\chi\J^M)+\chi\,\partial_M\J^M,
\end{equation}
so gauge invariance of the coupling (up to boundary terms) is equivalent to
\(\partial_M\J^M=0\).

\paragraph{Charge labels and reduced coupling.}
The updated charge ontology distinguishes the fixed defect-branch sign
\(\etaQ=\pm1\) from the positive microscopic scale \(\estar>0\), with
\(\qstar\equiv\etaQ\estar\). The canonically normalized brane coupling is then
\(\qeff=\qstar/\sqrt{\Zint}=\etaQ\,\eeff\), where
\(\eeff\equiv\estar/\sqrt{\Zint}\). Thus \(\qstar\) labels the microscopic defect
branch, while \(\qeff\) measures the observable brane coupling after zero-mode
normalization; neither sign nor magnitude is defined here by circulation or throat
breathing.

\paragraph{Gauge fixing.}
We include a Lorenz-type gauge-fixing term
\(-(\partial\!\cdot\!\A)^2/(2\xigf\muo)\) with
\(\partial\!\cdot\!\A\equiv\partial_M\A^M\). This term explicitly breaks gauge
invariance, as expected, but renders the operator invertible and makes the
divergence identity and component checks unambiguous.

\subsection{Lorentz transformations: covector vs vector}

A persistent source of apparent ``Lorentz non-invariance'' in component checks is
using the wrong transformation type for \(\A_\mu\). In this paper:

\begin{itemize}
\item \(\A_\mu\) transforms as a \emph{covector} (one-form):
\(\A'_\mu(x')=(\Lambda^{-1})^\nu{}_\mu\,\A_\nu(x)\).
\item \(\J^\mu\) transforms as a \emph{vector}:
\(\J'^\mu(x')=\Lambda^\mu{}_\nu\,\J^\nu(x)\).
\end{itemize}

With these assignments, \(\A_\mu\J^\mu\) is a Lorentz scalar and
\(\F_{\mu\nu}\F^{\mu\nu}\) is invariant. Treating \(\A_\mu\) as a vector instead
will generally produce spurious sign/index mismatches in invariance checks.

\subsection{Brane reduction: what must vanish, and why}

The controlled brane Maxwell limit relies on two closely related assumptions:

\begin{enumerate}
\item \textbf{Axial gauge and zero-mode dominance:} \(\A_w=0\) and
\(\partial_w\A_\mu\simeq 0\) over the localization region so that
\(\F^{w\nu}=-\partial^w\A^\nu\simeq 0\).
\item \textbf{No transverse source flux:} \(\J^w=0\), consistent with the reduced
\(N=w\) equation \cref{eq:eom_w}.
\end{enumerate}

If \(\J^w\neq 0\) or if significant \(w\)-dependence is present, the brane theory
departs from pure Maxwell in a way controlled by the KK tower, as discussed in
\cref{sec:discussion}. These conditions therefore define a controlled far-field
sector: they suppress the mixed core channels \((A_w,J^w,F_{\mu w})\) without
implying that such structure is absent from the microscopic defect ontology.

\subsection{Implementation pitfalls in symbolic workflows (WL)}

The accompanying WL suite is structured to avoid several common pitfalls that can
make correct physics look incorrect:

\paragraph{(i) \texttt{EulerEquations} returns \texttt{Equal[\dots]}.}
In WL, \texttt{EulerEquations} often returns equations as \texttt{lhs == rhs}.
When checking equality against a target expression, it is safer to convert to a
single expression via \(\texttt{lhs - rhs}\) before simplifying. Directly
simplifying a difference of \texttt{Equal} objects can yield misleading
nonzero residuals.

\paragraph{(ii) Protected iterators and malformed replacement rules.}
Using symbolic iterators in \texttt{Table} (e.g.\ \(\{a,b,c\}\) as if they were
numeric ranges) can trigger protected-symbol errors. Likewise, constructing
replacement rule lists indirectly (e.g.\ via nested \texttt{Join/Table} patterns)
can silently produce malformed rule sets. The suite avoids these issues by using
explicit index lists and clean, flat replacement-rule lists.

\paragraph{(iii) Symmetry checks require correct index raising/lowering.}
Sign mismatches often arise from mixing covariant and contravariant components
without explicitly applying \(\eta_{MN}\) or \(\eta^{MN}\). The suite always
computes \(\A^M=\eta^{MN}\A_N\) and \(\F^{MN}=\eta^{MP}\eta^{NQ}\F_{PQ}\) before
forming divergences or contractions.

\paragraph{(iv) Noether identity sign convention.}
The off-shell Noether identity used in \cref{eq:noether_identity} has a definite
relative sign between \(\partial_\mu j^\mu\) and the Euler--Lagrange combination
\((\psi^*\mathrm{EL}_{\psi^*}-\psi\,\mathrm{EL}_\psi)\). Using the opposite sign
will produce an apparent ``failure'' even if the extracted current is correct.
The suite implements the identity in the form written in the main text.

\paragraph{(v) KK norms for symbolic mode number.}
Direct symbolic integration of Hermite norms for generic \(n\) can return
\texttt{Indeterminate} in WL. The suite validates the closed form by spot-checking
integer \(n\) and then uses the validated closed-form expression for general
statements about \(c_n\) and the Yukawa tower.
% ============================================================

\end{document}
